% Tabla A3-T2. Generada automáticamente; no editar a mano.
% Requiere \usepackage{booktabs} (y graphicx si la tabla es ancha).
\begin{table}[htbp]
    \centering
    \caption{Análisis III: resumen de la muestra de sistemas multiplanetarios.}
    \label{tab:A3-T2}
    \small
    \resizebox{\textwidth}{!}{%
    \begin{tabular}{lp{0.40\textwidth}rr}
    \toprule
     & Característica & Valor & \% \\
    \midrule
    Muestra & Sistemas & 615 & -- \\
    Muestra & Planetas en esos sistemas & 1\,513 & -- \\
    Muestra & Sistemas gapped & 85 & 13.8 \\
    Medianas & Mediana de la razón de períodos adyacentes & 2.343 & -- \\
    Medianas & Mediana del radio planetario mediano [$R_\oplus$] & 2.045 & -- \\
    Medianas & Mediana de la masa estelar [$M_\odot$] & 0.885 & -- \\
    Medianas & Mediana de {[Fe/H]} [dex] & -0.006 & -- \\
    Multiplicidad & N = 2 & 415 & 67.5 \\
    Multiplicidad & N = 3 & 137 & 22.3 \\
    Multiplicidad & N = 4 & 47 & 7.6 \\
    Multiplicidad & N $\geq$ 5 & 16 & 2.6 \\
    Origen observacional & Kepler & 389 & 63.3 \\
    Origen observacional & TESS & 95 & 15.4 \\
    Origen observacional & solo transito & 67 & 10.9 \\
    Origen observacional & velocidad radial & 64 & 10.4 \\
    Radios & Sistemas con todos sus planetas observados en tránsito & 542 & 88.1 \\
    \bottomrule
    \end{tabular}
    }
    \par\smallskip
    {\footnotesize\textit{Nota:} Origen observacional: Kepler o TESS si alguna instalación de descubrimiento lo es; el resto, según el método de detección.}
\end{table}
